\documentclass[12pt,a4paper]{article}
\synctex=1
\usepackage[utf8]{inputenc}
\usepackage[margin=1cm]{geometry}
\usepackage{graphicx}
%\usepackage{verbatim}
\usepackage{listings}
\usepackage{textcomp}
\usepackage{courier}
\usepackage{libertine}
\usepackage{pgfornament}
\usepackage{eso-pic}
\usepackage[hangul]{kotex}
\linespread{1.3}

\title{
	\centering
	\pgfornament[width=12cm,color=teal]{84}\\
	\vspace{1cm}
	\fontsize{50}{50} \selectfont {TimeMachine Editor\\ with Java}\\
		\pgfornament[width=12cm,color=teal]{88}\\
	\vfill}
\author{
	\LARGE
	\begin{tabular}{rl}
		\hline
		학번 : & 2016110056\\ 
		학과 : & 불교학부 \\
		이름 : & 박승원\\
		날짜 : & \today\\
		\hline
	\end{tabular}\vspace{2cm}
	\\
\includegraphics[width=0.5\textwidth]{logo.jpg}
	}
\date{}


\begin{document}
%\maketitle
\pagenumbering{gobble}
\noindent
\lstset{language=java, columns=flexible, tabsize=4, frame=shadowbox, showstringspaces=false, breaklines=true, upquote=true, basicstyle=\normalsize}

\section{Functionality}
\begin{itemize}
	\item basic editing
	\item show line number 
	\item open \& write file
	\item search string with case sensitive option \textrightarrow show the line in Italic font
	\item \large\textbf revert to any moment of the document
\end{itemize}

\section{Algorithm}
In vector graphic images, we save all the drawing procedures so that we can draw any scale of the designated image.
Likewise this editor saves all the keystrokes so that we can go back to any moment of the document by repeating the keystrokes from the first.

So this editor saves two files. One is plain text file and the other is integer values of keystrokes and revert-records which is expressed in negative integer.

\begin{lstlisting}
	public void backtothefuture(int year) {
		int n = keyrecord.size();
		record = false;
		clear();
		for(int i=0; i<year; i++) {
			int el = keyrecord.elementAt(i);
			if(el < 0) backtothefuture(i + el);
			else input(el);
			if(n != keyrecord.size()) System.out.println("error"); 
		}
	}
\end{lstlisting}

In upper program is shown the recursive revert procedure. 
When el is negative, it goes that much keystrokes back. 
\section{Implementation}
I had to avoid using JTextArea to record all the keystrokes.
I used JComponent as textArea.
Cursor is shown with vertical line.

\newpage
Basic View

\includegraphics[width=0.8\textwidth]{1.png}

\section{Usage}
Usage is very intuitive that no explanation is needed except reverting document.

There is a slider at the bottom of the window.
Right most part of the slider is present and left most is past.
Slider automatically correct the minimum and maximum of the values according to keystrokes entered.
If it is moved, text shown will be changed to corresponding moment.
When you decided a moment you want to go back, click on the button near it.
And continue writing.

All the procedure is recorded including the part that is canceled by going back.

\section{Benefit}
If we use this editor with version control system, we can greatly reduce needed commit. Only at the branching moment will we need the commit.

And in normal situations like writing a email, this editor can be used to ease the editing process. Especially when you are going to edit it a lot.
\end{document}
